% --------------------------------------------------------------------------
% MACROS AUXILIARES
% --------------------------------------------------------------------------
\newcommand{\code}[1]{\texttt{\detokenize{#1}}}
\newcommand{\pendente}[1]{\textbf{\textcolor{red!80!black}{[PENDENTE: #1]}}}
\newcommand{\ph}[1]{\textcolor{red!60!black}{\textit{[#1]}}}
\newcommand{\rail}{\num{8.5}\,$\times$\,\num{8.5}~\si{\milli\metre}}
\newcommand{\figplaceholder}{\fbox{\rule{0pt}{5cm}\rule{0.8\textwidth}{0pt}}}
\newcommand{\Grmsbase}{G_{\text{rms, base}}}
\newcommand{\Grmspay}{G_{\text{rms, payload}}}

% --------------------------------------------------------------------------
% LAYOUT ABNT: margens (sup./esq. 3 cm; inf./dir. 2 cm), 1,5 entrelinhas,
% recuo de parágrafo de 1,25 cm
% --------------------------------------------------------------------------
\setlrmarginsandblock{3cm}{2cm}{*}
\setulmarginsandblock{3cm}{2cm}{*}
\checkandfixthelayout
\OnehalfSpacing
\setlength{\parindent}{1.25cm}
\setlength{\parskip}{0pt}
\setcounter{secnumdepth}{5}
\setcounter{tocdepth}{3}
\renewcommand{\orientadorname}{Orientadora:}
\sloppy

% Capa Institucional UESC conforme NBR 14724
\renewcommand{\imprimircapa}{%
  \begin{capa}%
    \center
    {\ABNTEXchapterfont\bfseries\large UNIVERSIDADE ESTADUAL DE SANTA CRUZ -- UESC\par}
    \vspace*{0.25cm}
    {\ABNTEXchapterfont\bfseries\large DEPARTAMENTO DE CIÊNCIAS EXATAS E TECNOLÓGICAS\par}
    \vspace*{0.25cm}
    {\ABNTEXchapterfont\bfseries\large COLEGIADO DE ENGENHARIA MECÂNICA\par}
    \vspace*{3.5cm}
    {\ABNTEXchapterfont\bfseries\large\imprimirautor\par}
    \vfill
    \begin{center}
      \ABNTEXchapterfont\bfseries\LARGE\imprimirtitulo
    \end{center}
    \vfill
    {\large\imprimirlocal\par}
    \vspace*{0.2cm}
    {\large\imprimirdata\par}
    \vspace*{1cm}
  \end{capa}
}

% Lista de referências: sem rótulos, espaçamento simples (NBR 6023)
\makeatletter
\renewcommand{\@biblabel}[1]{}
\makeatother

% --------------------------------------------------------------------------
% HIPERLIGAÇÕES DISCRETAS
% --------------------------------------------------------------------------
\hypersetup{
  colorlinks   = true,
  linkcolor    = black,
  citecolor    = black,
  urlcolor     = black,
  pdftitle     = {CubeSat em Baixa Órbita: Otimização de Volume e Geometria para Maximização de Desempenho Estrutural e Funcional},
  pdfauthor    = {Gabriel Marques de Andrade},
  pdfkeywords  = {CubeSat, Metamateriais Auxéticos, AlSi10Mg, NASA GEVS},
  bookmarksdepth = 3
}

% --------------------------------------------------------------------------
% DADOS DE IDENTIFICAÇÃO
% --------------------------------------------------------------------------
\titulo{CubeSat em Baixa Órbita: Otimização de Volume e Geometria para Maximização de Desempenho Estrutural e Funcional}
\autor{Gabriel Marques de Andrade}
\local{Ilhéus -- BA}
\data{2026}
\orientador{Prof.ª Dr.ª Nila Cecília de Faria Lopes Medeiros}
\instituicao{%
  UNIVERSIDADE ESTADUAL DE SANTA CRUZ -- UESC
  \par DEPARTAMENTO DE CIÊNCIAS EXATAS E TECNOLÓGICAS
  \par COLEGIADO DE ENGENHARIA MECÂNICA}
\tipotrabalho{Trabalho de Conclusão de Curso}
\preambulo{Trabalho de Conclusão de Curso apresentado ao Colegiado do Curso de Engenharia Mecânica da Universidade Estadual de Santa Cruz (UESC), como requisito parcial para a obtenção do título de Bacharel em Engenharia Mecânica.}

